\documentclass[12pt]{article}

\usepackage[margin=1in]{geometry}
\usepackage{amsmath,amssymb,amsthm,amsfonts}
\usepackage{mathtools}
\usepackage{natbib}
\usepackage{hyperref}
\usepackage{bm}

\hypersetup{
	colorlinks=true,
	linkcolor=blue,
	citecolor=blue,
	urlcolor=blue
}

% Theorem environments
\newtheorem{theorem}{Theorem}[section]
\newtheorem{proposition}[theorem]{Proposition}
\newtheorem{lemma}[theorem]{Lemma}
\newtheorem{corollary}[theorem]{Corollary}
\newtheorem{definition}[theorem]{Definition}
\newtheorem{remark}[theorem]{Remark}

\numberwithin{equation}{section}

\begin{document}
	
	\title{Minimum Complexity Estimation:\\
		Oracle Inequalities, Minimax Bounds and Asymptotic Theory}
	\author{J.~Nembe}
	\date{\today}
	\maketitle
	
	\begin{abstract}
		We study a general class of \emph{minimum complexity estimators} (MCE),
		where the estimator is defined as a penalized likelihood minimizer with
		a complexity penalty interpreted as a description length or code length.
		Under a Kraft-type inequality on the penalties, we derive nonasymptotic
		oracle inequalities for the Hellinger risk (the index of resolvability),
		minimax lower bounds that match the oracle (index-of-resolvability)
		rate up to constants, as well as consistency and
		asymptotic normality under standard regularity conditions. We further
		pin down the precise minimax status of the procedure: in the parametric
		regime the estimator is minimax up to constants, whereas on a
		nonparametric sieve its description-length penalty inflates the rate by a
		logarithmic factor; we show this factor is a coding artefact, removed
		exactly---for known smoothness and adaptively---by a penalized-projection
		refinement that pays the model's metric complexity rather than the
		per-coefficient quantization cost. The proofs
		combine information-theoretic arguments with empirical process techniques
		and show that the minimum complexity principle yields theoretically sound
		procedures that adapt to the unknown model complexity.
	\end{abstract}
	
	\tableofcontents
	
\section{Introduction}

Model selection and prediction under general nonparametric or
semiparametric families are among the central problems of modern
statistical theory.  Since the seminal works of Akaike, Schwarz,
Barron, Birgé and Massart, a vast literature has shown that
penalized likelihood procedures achieve near–optimal risk bounds under
appropriate choices of penalties.  More recent contributions, such as
van de Geer (2000), have refined these results through local
entropy methods and sharp empirical process techniques.

Despite this rich body of work, a persistent limitation remains:
\emph{classical oracle inequalities are fundamentally tied to the
	geometry of the likelihood and to the concentration of the empirical
	process, rather than to the complexity of the underlying models
	themselves}.  In particular, classical approaches rely on:
\begin{itemize}
	\item uniform deviation bounds over models indexed by metric
	complexities (entropy, bracketing, VC dimension, etc.);
	\item parametric or local asymptotic arguments for well–behaved
	likelihoods;
	\item explicitly chosen penalties tuned to the geometry of the model.
\end{itemize}

In contrast, this paper develops a general theory of
\emph{minimum–complexity estimation (MCE)}, in which the selection
criterion is driven purely by \emph{code length}––that is,
by the amount of information needed to specify the model itself.
The central objective is to obtain a unified framework where
penalized likelihood, universal coding, and KL–based risk minimization
all arise from a single principle:
\[
(\widehat{m},\widehat{\theta})
\in \arg\min_{m,\theta}
\Big\{-\log L_n(f_{m,\theta}) + L(m)\Big\},
\]
with $L(m)$ satisfying Kraft’s inequality.  
This formulation yields an estimator that is conceptually simple,
computationally direct, and theoretically powerful.

\paragraph{Main contributions.}
We show that the minimum–complexity estimator satisfies a
\emph{non-asymptotic, model–wise oracle inequality} for the squared
Hellinger risk $h^2(P^{\star},\widehat{f})$ (the Barron--Cover index of
resolvability)
\[
\mathbb{E}\Big[
h^2\!\big(P^{\star},\widehat{f}\big)
\Big]
\le
C\!\left(
\inf_{m\in\mathcal{I}}
\Big\{
\inf_{\theta\in\Theta_m}
\mathrm{KL}(P^{\star}\,\|\,f_{m,\theta})
+ \frac{L(m)}{n}
\Big\}
\right),
\]
valid \emph{without metric-entropy, curvature, margin, or sub–model
regularity conditions beyond Kraft's inequality} (for a countable family;
a continuous $\Theta_m$ requires augmenting $L(m)$ by the code length of a
discretization of $\theta$). The Hellinger risk is essential here: the
analogous bound for the \emph{Kullback--Leibler} risk does not hold without
a further boundedness condition on the log-likelihood ratios, since the
penalized maximizer may select a density assigning very small mass where
$P^\star$ does not, making $\mathrm{KL}(P^\star\,\|\,\widehat f)$ infinite
while the resolvability term stays finite.
This is in sharp contrast with the classical theory, where metric-complexity
assumptions play an essential role in controlling the empirical
process.

The proof rests on a single Kraft/partition-function bound for the
Hellinger affinity, rather than bounding the supremum of
the empirical process.
As a result, no geometry of the parameter space is required beyond
the Kraft condition.  This provides a form of universality:
\emph{any countable model family endowed with prefix code lengths
	automatically satisfies the oracle inequality}.  

We further show that the MCE enjoys:
\begin{itemize}
	\item a minimax lower bound showing that the rate is
	essentially optimal for the entire model family;
	\item strong consistency and asymptotic normality under mild
	identifiability and regularity, with no need for local
	bracketing entropy;
	\item a structural interpretation in terms of universal coding,
	connecting estimation to information–theoretic principles.
\end{itemize}

\paragraph{Differences with Barron–Birgé–Massart and van de Geer.}
Although this work is philosophically aligned with the universal coding
approach of Barron and with the penalized likelihood program developed
by Birgé and Massart, the present results differ in three fundamental
ways:

\begin{enumerate}
	\item \textbf{Elimination of geometric complexity assumptions.}
	Classical penalties such as AIC, BIC, or MDL rely on model
	dimensions, local entropy conditions, or the curvature of the
	log–likelihood.  
	Our bound requires \emph{no such assumptions}: Kraft’s inequality
	alone suffices, and the empirical process does not need to be
	controlled on model–dependent neighborhoods.
	
	\item \textbf{Uniformity over arbitrarily large model families.}
	In the Barron–Birgé–Massart framework, the complexity of the
	model family must be controlled a priori
	(through entropy or bracketing integrals).  
	Here, the result applies immediately to \emph{any} countable set
	of models with valid code lengths.
	
	\item \textbf{A single structural principle.}
	Van de Geer’s theory is based on local empirical process methods,
	margin conditions, and intricate contraction arguments.
	The present approach derives everything from a canonical
	likelihood–complexity decomposition and a coding–theoretic argument.
	The resulting bound is therefore both simpler and fundamentally different:
	\emph{the estimator is optimal because its description length is optimal}.
\end{enumerate}

Thus, the contribution of this article is not merely to provide another
oracle inequality, but to show that oracle optimality can be achieved
in full generality from a single information–theoretic principle, with
no geometric assumptions on the underlying models.
This aligns statistical estimation with a universal coding perspective,
offering a clean, minimalistic and powerful foundation for general
model selection.

	
	\section{Framework and Minimum Complexity Estimator}
	\label{sec:framework}
	
	Let $(\mathcal{X},\mathcal{A})$ be a measurable space endowed with a
	$\sigma$-finite dominating measure $\mu$. The true distribution
	$P^{\star}$ is assumed to admit a density $f^{\star}$ with respect to
	$\mu$.
	
	We consider a countable index set $\mathcal{I}$ and, for each
	$m\in\mathcal{I}$, a model $\mathcal{M}_m$ consisting of densities
	$f_{m,\theta}$ with $\theta\in\Theta_m$:
	\[
	\mathcal{M}_m = \{f_{m,\theta}: \theta\in\Theta_m\}.
	\]
	We let $\mathcal{M} = \bigcup_{m\in\mathcal{I}}\mathcal{M}_m$ denote
	the full model family.
	
	Given an i.i.d.\ sample $(X_1,\ldots,X_n)$ distributed according to
	$P^{\star}$, the log--likelihood of $f_{m,\theta}$ is
	\[
	\ell_n(m,\theta)
	:= \sum_{i=1}^n \log f_{m,\theta}(X_i).
	\]
	
	\subsection{Complexity penalties and Kraft inequality}
	
	To encode the model index $m$ we consider a \emph{code length} function
	$L:\mathcal{I}\to[0,\infty)$ satisfying the Kraft inequality
	\begin{equation}
		\label{eq:kraft-main}
		\sum_{m\in\mathcal{I}} 2^{-L(m)} \le 1.
	\end{equation}
	This ensures the existence of a prefix code with code lengths $L(m)$,
	and provides a natural information-theoretic interpretation of the
	penalty.
	
	Typical choices include, for example,
	\[
	L(m) \asymp \dim(\Theta_m)\,\log n
	\]
	in parametric or sieve-type settings, though our results do not require
	any specific form as long as \eqref{eq:kraft-main} holds.
	
	\subsection{Definition of the minimum complexity estimator}
	
	The minimum complexity estimator (MCE) is defined as any measurable
	selection
	\begin{equation}
		\label{eq:mce-def-main}
		(\widehat{m},\widehat{\theta})
		\in \arg\min_{m\in\mathcal{I},\ \theta\in\Theta_m}
		\Big\{
		-\ell_n(m,\theta)
		+ L(m)
		\Big\}.
	\end{equation}
	We then write
	\[
	\widehat{f} := f_{\widehat{m},\widehat{\theta}}
	\]
	for the estimated density.
	
	\begin{remark}
		The penalty is placed solely on the model index $m$. In many practical
		situations one can refine this to include a parametric complexity term
		depending on $\theta$ but, for clarity, we focus on the purely
		model-based penalty here.
	\end{remark}
	
	\section{Main Results}
	\label{sec:main-results}
	
	In this section we state the main theoretical results concerning the
	risk, minimax properties, and asymptotic behavior of the MCE.
	
	\subsection{Oracle inequality}
	
	Our first theorem is a nonasymptotic oracle inequality for the
	Kullback--Leibler risk.
	
	\begin{theorem}[Oracle inequality for the MCE]
		\label{thm:oracle}
		Assume that $(L(m))_{m\in\mathcal{I}}$ satisfies the Kraft inequality
		\eqref{eq:kraft-main}. Let $\widehat{f}$ be any MCE defined by
		\eqref{eq:mce-def-main}. Then there exists a universal constant $C>0$
		such that, for the squared Hellinger distance
		$h^2(P,Q)=\tfrac12\int(\sqrt{dP}-\sqrt{dQ})^2$,
		\begin{equation}
			\label{eq:oracle-statement}
			\mathbb{E}\Big[
			h^2\!\big(P^{\star},\widehat{f}\big)
			\Big]
			\le C \inf_{m\in\mathcal{I}}
			\left\{
			\inf_{\theta\in\Theta_m}
			\mathrm{KL}\!\big(P^{\star}\,\|\,f_{m,\theta}\big)
			+ \frac{L(m)}{n}
			\right\}.
		\end{equation}
		(For a continuous $\Theta_m$, the penalty $L(m)$ must be augmented by the
		code length of a discretization of $\theta$; see the proof.)
	\end{theorem}
	
	\begin{remark}
		The inequality \eqref{eq:oracle-statement} shows that the MCE
		automatically adapts to the best model in the collection, up to a
		complexity term $L(m)/n$ that reflects the description length of the
		model index. In particular, if $P^{\star}$ is well approximated by a
		low-complexity model, the risk of the MCE will be correspondingly low.
	\end{remark}
	
	\subsection{Minimax lower bound}
	
	We complement the oracle inequality with a minimax lower bound showing
	that the oracle rate cannot be improved in a distribution-free sense
	(up to constants).
	
	\begin{theorem}[Minimax lower bound]
		\label{thm:minimax-lower}
		Let $\mathcal{F}$ be a model class containing a finite subset
		$\{f_0,\ldots,f_M\}$ such that
		\begin{enumerate}
			\item $\mathrm{KL}(f_j\,\|\,f_k) \le \alpha$ for all $j\ne k$;
			\item $\mathrm{KL}(f_j\,\|\,f_0) \ge \beta$ for all $j\ge 1$,
		\end{enumerate}
		for some $\alpha,\beta>0$. Then, for any estimator $\widehat{f}$ based
		on $n$ observations,
		\begin{equation}
			\label{eq:minimax-statement}
			\sup_{f\in\mathcal{F}}
			\mathbb{E}_f\big[\mathrm{KL}(f\,\|\,\widehat{f})\big]
			\ge c\,\beta\,
			\left(1 - \frac{n\alpha + \log 2}{\log M}\right),
		\end{equation}
		for a universal constant $c>0$.
	\end{theorem}
	
	\begin{remark}
		Theorem~\ref{thm:minimax-lower} is a standard Fano-type statement. In
		concrete settings, one can construct finite subsets $\{f_j\}$ with
		pairwise Kullback--Leibler distances controlled as in (i)–(ii) and thereby
		obtain explicit lower bounds in terms of $n$, the dimension of the model,
		or its metric entropy. The result shows that the rate appearing in the
		oracle inequality is essentially optimal. The precise gap between this
		oracle (resolvability) rate and the \emph{exact} minimax rate on
		nonparametric sieves---a logarithmic factor of purely coding origin---is
		analyzed in Section~\ref{sec:log-gap}.
	\end{remark}

	\subsection{From near-minimax to exact: coding cost versus metric complexity}
	\label{sec:log-gap}

	The oracle inequality \eqref{eq:oracle-statement} is sharp \emph{as a bound on
	the index of resolvability}. When instantiated on a nonparametric sieve,
	however, the description-length penalty $L(m)$ carries a logarithmic factor
	that is \emph{absent} from the minimax rate. We isolate this factor, show it is
	a coding artefact, and remove it.

	Consider the canonical sieve: $\mathcal{M}_k$ is a linear space of dimension
	$D_k\asymp 2^{kd}$ (e.g.\ piecewise polynomials of degree $<\lceil s\rceil$ at
	resolution $2^{-k}$) approximating a bounded density $f^{\star}$ with
	$0<c_0\le f^{\star}\le C_1<\infty$ and Hölder smoothness $s$ on a
	$d$-dimensional domain, so that
	$\inf_{g\in\mathcal{M}_k}\mathrm{KL}(P^{\star}\,\|\,g)\le C\,L^2\,2^{-2ks}$.

	\begin{proposition}[The MCE rate carries a logarithmic factor]
		\label{prop:mce-log}
		To obtain a countable family the coefficients of $\mathcal{M}_k$ must be
		quantized on a grid of spacing $\delta_k$; with $\delta_k\asymp 2^{-ks}$
		this costs $L(k)\asymp D_k\log_2(1/\delta_k)\asymp D_k\,ks\lesssim
		D_k\log n$ bits. The oracle inequality \eqref{eq:oracle-statement} then
		gives
		\[
		\mathbb{E}\big[h^2(P^{\star},\widehat f)\big]
		\le C\inf_k\Big\{L^2 2^{-2ks}+\frac{D_k\log n}{n}\Big\}
		\le C'\Big(\frac{\log n}{n}\Big)^{2s/(2s+d)},
		\]
		i.e.\ the minimax rate $n^{-2s/(2s+d)}$ \emph{up to a logarithmic factor}.
	\end{proposition}

	The $\log n$ is thus a \emph{coding artefact}: it is the per-coefficient cost
	$\log(1/\delta_k)$ of describing a continuum by a prefix code, not a feature of
	the estimation problem. An estimator that fits the coefficients of
	$\mathcal{M}_k$ \emph{without} quantizing them pays instead the metric
	complexity $D_k/n$ of the model and attains the exact rate.

	\begin{theorem}[Exact minimax rate via penalized projection]
		\label{thm:exact-proj}
		Let $\widehat g_k=\Pi_k P_n$ be the orthogonal $L^2(\mu)$-projection of the
		empirical measure onto $\mathcal{M}_k$, and let $\widehat f$ be its
		truncation to $[c_0/2,2C_1]$, renormalized to a density.
		\begin{enumerate}
			\item \emph{(Known $s$.)} The single resolution $D_{k_n}\asymp
			n^{d/(2s+d)}$ gives
			$\mathbb{E}\,h^2(P^{\star},\widehat f_{k_n})\le C\big(2^{-2k_n s}
			+D_{k_n}/n\big)\le C'\,n^{-2s/(2s+d)}$.
			\item \emph{(Adaptive.)} The penalized selection $\widehat k\in
			\arg\min_k\{-\|\widehat g_k\|_2^2+\kappa D_k/n\}$ attains
			$\mathbb{E}\,h^2(P^{\star},\widehat f)\le C\inf_k\{L^2 2^{-2ks}+D_k/n\}
			\le C'\,n^{-2s/(2s+d)}$, without knowledge of $s$.
		\end{enumerate}
		In both cases there is \emph{no} logarithmic factor.
	\end{theorem}

	\begin{proof}[Proof sketch]
		(i) Orthogonality of the projection splits
		$\mathbb{E}\|\widehat g_k-f^{\star}\|_2^2$ into the approximation term
		$\|f^{\star}-\Pi_k f^{\star}\|_2^2\le C L^2 2^{-2ks}$ and the variance
		$n^{-1}\sum_{\ell\le D_k}\mathrm{Var}\big(e_\ell(X)\big)\le C_1 D_k/n$,
		where $(e_\ell)$ is an $L^2(\mu)$-orthonormal basis of $\mathcal{M}_k$ and
		$\sum_\ell\|e_\ell\|_2^2=D_k$. The truncation is an $L^2$-contraction toward
		an interval containing $f^{\star}$, and under the two-sided bound
		$f^{\star},\widehat f\in[c_0/2,2C_1]$ one has $h^2\lesssim\|\cdot\|_2^2$;
		balancing $2^{-2k_n s}\asymp D_{k_n}/n$ gives $n^{-2s/(2s+d)}$.
		(ii) The criterion is the least-squares contrast
		$\gamma_n(g)=\|g\|_2^2-2P_n g$, minimized over $\mathcal{M}_k$ by
		$\widehat g_k$. The Birgé--Massart penalized-projection oracle inequality
		\citep{BirgeMassart1998,Massart2007} applies: since the dimensions
		$D_k\asymp 2^{kd}$ grow \emph{geometrically}, the weights $x_k\asymp D_k$
		are summable ($\sum_k e^{-D_k}<\infty$), so the penalty
		$\mathrm{pen}(k)\asymp D_k/n$---proportional to the metric complexity, with
		no extra $\log$---is admissible and yields
		$\mathbb{E}\|\widehat g_{\widehat k}-f^{\star}\|_2^2\le
		C\inf_k\{\|f^{\star}-\Pi_k f^{\star}\|_2^2+D_k/n\}$. Inserting the
		approximation bound gives the exact adaptive rate.
	\end{proof}

	\begin{remark}[Honest minimax status of the MCE]
		\label{rem:honest-mce}
		Theorem~\ref{thm:oracle} shows the MCE is oracle-optimal \emph{for the index
		of resolvability}, and Theorem~\ref{thm:minimax-lower} shows this index is
		optimal up to constants. But on a nonparametric sieve the resolvability rate
		itself exceeds the minimax rate by the coding factor $\log n$
		(Proposition~\ref{prop:mce-log}); the MCE is then \emph{near-minimax}
		(optimal up to a logarithmic factor), not minimax up to constants. The exact
		rate is recovered by the penalized-projection estimator of
		Theorem~\ref{thm:exact-proj}, which differs from the MCE precisely by
		replacing the description-length penalty $D_k\log n$ with the
		metric-complexity penalty $D_k$. In the \emph{parametric} regime
		(finite-dimensional $\Theta_m$, bounded $L(m)$) no such factor arises and the
		MCE is minimax up to constants. The same dichotomy, instantiated on a
		$G$-invariant density on a quotient $E/G$ of effective dimension
		$d_{\mathrm{eff}}=\dim(E/G)$, underlies the orbital minimum-complexity
		estimator: there too the description-length code is near-minimax and the
		projection estimator is exact.
	\end{remark}

	\subsection{Consistency and asymptotic normality}
	
	We now state the main asymptotic results for the MCE.
	
	\begin{theorem}[Strong consistency]
		\label{thm:consistency}
		Assume that:
		\begin{enumerate}
			\item The true distribution $P^{\star}$ belongs to
			$\bigcup_{m\in\mathcal{I}}\mathcal{M}_m$ and is uniquely represented
			as $f_{m^{\star},\theta^{\star}}$;
			\item For each $m$, the parameter space $\Theta_m$ is compact and
			the map $\theta\mapsto f_{m,\theta}$ is continuous in
			$\mathrm{KL}$;
			\item The penalty satisfies $L(m) = o(n)$ as $n\to\infty$ for each
			fixed $m$ and grows sufficiently fast in $m$ to exclude overly
			complex models.
		\end{enumerate}
		Then
		\[
		\widehat{m}\to m^{\star},\qquad
		\widehat{\theta}\to\theta^{\star}
		\]
		almost surely as $n\to\infty$.
	\end{theorem}
	
	Under additional smoothness assumptions we obtain asymptotic normality.
	
	\begin{theorem}[Asymptotic normality]
		\label{thm:asymp-normal}
		In addition to the hypotheses of Theorem~\ref{thm:consistency}, assume
		that the model $\mathcal{M}_{m^{\star}}$ is regular in the sense of
		classical likelihood theory, with Fisher information matrix
		$I(\theta^{\star})$ positive definite. Then
		\[
		\sqrt{n}\,(\widehat{\theta}-\theta^{\star})
		\xrightarrow{d}
		\mathcal{N}\big(0, I(\theta^{\star})^{-1}\big).
		\]
	\end{theorem}
	
	\section{Nonasymptotic Tools and Empirical Processes}
	\label{sec:nonasymptotic}
	
	We briefly sketch the main probabilistic ingredients used in the proofs.
	
	First, we exploit a basic decomposition of the Kullback--Leibler risk in
	terms of likelihood ratios (Lemma~\ref{lem:KL-decomp} in
	Section~\ref{sec:proofs}) combined with the optimality condition
	\eqref{eq:mce-def-main} to compare $\widehat{f}$ to any fixed
	$f_{m,\theta}$. This yields an inequality of the form
	\[
	\mathrm{KL}(P^{\star}\,\|\,\widehat{f})
	\lesssim \mathrm{KL}(P^{\star}\,\|\,f_{m,\theta})
	+ \frac{L(\widehat{m})-L(m)}{n}
	+ \text{empirical process fluctuations},
	\]
	which is then controlled via exponential inequalities for empirical
	processes indexed by log--likelihood differences.
	
	Second, we use simple Bernstein-type inequalities and, when necessary,
	entropy bounds for function classes to control the suprema of empirical
	processes (Lemma~\ref{lem:exp-ineq} and Lemmas~\ref{lem:covering}--\ref{lem:local-entropy}
	in Section~\ref{sec:proofs}). This allows us to obtain high-probability
	versions of the oracle inequality.
	
	Third, for the minimax lower bound we rely on a classical Fano argument
	(see, e.g., \citealp{Tsybakov2009}). The key idea is to embed a finite
	subset of well-separated models into $\mathcal{M}$ and relate the
	estimation problem to a multi-hypothesis testing problem.
	
	Finally, the consistency and asymptotic normality results follow the
	standard M-estimation scheme developed in \cite{vdVart1998}, after
	verifying that the penalized criterion behaves like the negative
	log--likelihood in a neighborhood of $(m^{\star},\theta^{\star})$ and
	that the penalty does not perturb the local quadratic approximation.
	
	\section{Simulation Illustration (Optional)}
	\label{sec:simulations}
	
	In a full version of the paper, this section would contain simulation
	experiments comparing the MCE with other model selection criteria such
	as AIC, BIC or cross-validation in various parametric and semi-parametric
	settings. The aim would be to illustrate:
	\begin{itemize}
		\item Adaptation of the MCE to the true model complexity;
		\item Robustness of the procedure under mild misspecification;
		\item Empirical verification of the oracle-type behavior suggested
		by Theorem~\ref{thm:oracle}.
	\end{itemize}
	
	\section{Discussion}
	\label{sec:discussion}
	
	We have shown that the minimum complexity principle, when formalized
	as a penalized likelihood procedure with code-length penalties, enjoys
	strong statistical guarantees. The resulting estimator achieves
	oracle-like performance for the index of resolvability, is minimax optimal
	up to constants in the parametric regime---and up to a logarithmic factor on
	nonparametric sieves, a coding cost removable by the penalized-projection
	refinement of Section~\ref{sec:log-gap}---and behaves like a regular
	M-estimator within the true model.
	
	Several extensions are natural. One direction is to consider
	data-dependent or hierarchical penalties, possibly reflecting
	\emph{twin} structures or group actions in the underlying models.
	Another direction is to strengthen the nonasymptotic analysis using
	refined chaining bounds and local Rademacher complexities, which may
	yield sharper constants and allow for heavier-tailed observations.
	
	Finally, the information-theoretic viewpoint suggests a deeper
	connection between minimum complexity estimation, universal source
	coding and the geometry of model spaces---an aspect that is especially
	relevant in the broader Twin Spaces program.
	
	\section{Proofs}
	\label{sec:proofs}
	
	In this section we provide complete proofs of the main results stated in
	Sections~\ref{sec:framework}--\ref{sec:main-results}.  Throughout, we work
	on the same probability space and use the notation introduced in the body
	of the paper.  For concreteness, we denote by $P^{\star}$ the unknown
	data--generating distribution, by $(\mathcal{M}_m)_{m\in\mathcal{I}}$ the
	model family, and by $f_{m,\theta}$ the densities (or likelihoods) with
	respect to a dominating measure $\mu$.  The \emph{minimum complexity
		estimator} (MCE) is written $\widehat{f} = f_{\widehat{m},\widehat{\theta}}$.
	
	\subsection{Preliminary inequalities}
	
	We begin with two elementary lemmas that will be used repeatedly.
	
	\begin{lemma}[Empirical/expected log-likelihood decomposition]
		\label{lem:KL-decomp}
		Let $P^{\star}$ have density $f^{\star}$ and let $Q$ be any probability
		measure with density $q$ (both with respect to $\mu$). Set
		$h_Q:=\log(f^{\star}/q)$. Then
		\begin{equation}
			P^{\star} h_Q = \mathrm{KL}(P^{\star}\,\|\,Q),
			\qquad
			P_n h_Q = \mathrm{KL}(P^{\star}\,\|\,Q) + (P_n-P^{\star})h_Q,
		\end{equation}
		where $P_n=n^{-1}\sum_{i=1}^n\delta_{X_i}$. Thus the empirical
		log-likelihood ratio equals the Kullback--Leibler divergence plus a
		mean-zero empirical-process fluctuation $(P_n-P^{\star})h_Q$.
	\end{lemma}

	\begin{proof}
		By definition $P^{\star}h_Q=\int f^{\star}\log(f^{\star}/q)\,d\mu
		=\mathrm{KL}(P^{\star}\,\|\,Q)$. The second identity is the trivial
		decomposition $P_n h_Q = P^{\star}h_Q + (P_n-P^{\star})h_Q$, and
		$\mathbb{E}_{P^{\star}}[(P_n-P^{\star})h_Q]=0$.
	\end{proof}

	\begin{remark}
		It is precisely the fluctuation $(P_n-P^{\star})h_{\widehat f}$, evaluated
		at the \emph{data-dependent} $\widehat f$, that must be controlled: the
		in-sample average $P_n h_{\widehat f}$ is \emph{not} equal in expectation
		to $\mathbb{E}[\mathrm{KL}(P^{\star}\,\|\,\widehat f)]$, because $\widehat f$
		is selected using the same sample. The proof of Theorem~\ref{thm:oracle-proof}
		below handles this through a Kraft/partition-function bound on the
		Hellinger affinity rather than through such a (false) in-sample
		identity.
	\end{remark}
	
	\begin{lemma}[Exponential inequality for empirical processes]
		\label{lem:exp-ineq}
		Let $(X_i)_{i=1}^n$ be i.i.d.\ with law $P^{\star}$ and let $\mathcal{F}$
		be a countable class of measurable functions such that
		$\|f\|_{\infty} \le b$ and $P^{\star}(f^2)\le v^2$ for all $f\in\mathcal{F}$.
		Then, for every $t>0$,
		\begin{equation}
			\label{eq:bernstein}
			\mathbb{P}\!\left(
			\sup_{f\in\mathcal{F}}
			\big|(P_n-P^{\star})f\big|
			\ge t
			\right)
			\le 2\,|\mathcal{F}|\,
			\exp\!\left(
			-\frac{n t^2}{2(v^2 + b t/3)}
			\right),
		\end{equation}
		where $P_n = n^{-1}\sum_{i=1}^n \delta_{X_i}$ is the empirical measure.
	\end{lemma}
	
	\begin{proof}
		For each fixed $f$, Bernstein's inequality yields
		\[
		\mathbb{P}\big(|(P_n-P^{\star})f|\ge t\big)
		\le 2 \exp\!\left(
		-\frac{n t^2}{2(v^2 + b t/3)}
		\right).
		\]
		The result follows by a union bound over the finite set $\mathcal{F}$.
		For a countable class one can approximate it by an increasing sequence of
		finite subclasses and use monotone convergence.
	\end{proof}
	
	\subsection{Proof of the basic oracle inequality}
	
	We now give a detailed proof of the main non--asymptotic oracle inequality
	for the minimum complexity estimator.  For ease of reference we restate
	its conclusion in a slightly simplified form; the precise constants can
	be recovered by keeping track of the intermediate bounds.
	
	\begin{theorem}[Oracle inequality for the MCE]
		\label{thm:oracle-proof}
		Assume that the code--lengths $(L(m))_{m\in\mathcal{I}}$ satisfy the Kraft
		inequality
		\begin{equation}
			\label{eq:kraft}
			\sum_{m\in\mathcal{I}} 2^{-L(m)} \le 1,
		\end{equation}
		and that for each model $\mathcal{M}_m$ there exists a (possibly
		data--dependent) penalized likelihood estimator $\widehat{\theta}_m$
		such that
		\begin{equation}
			\label{eq:def-mce}
			(\widehat{m},\widehat{\theta})
			\in \arg\min_{m,\theta}
			\Big\{-\log L_n(f_{m,\theta}) + L(m)\Big\},
		\end{equation}
		where $L_n(f_{m,\theta})$ is the likelihood of the sample
		$(X_1,\ldots,X_n)$ under $f_{m,\theta}$. Assume moreover (for definiteness)
		that the family is countable, each model contributing a single density
		$f_m$; for a continuous $\Theta_m$ the code length $L(m)$ is augmented by
		that of a discretization of $\theta$.
		Then, for the squared Hellinger distance $h^2$,
		\begin{equation}
			\label{eq:oracle-ineq}
			\mathbb{E}\Big[
			h^2\!\big(P^{\star},\widehat{f}\big)
			\Big]
			\le \inf_{m\in\mathcal{I}}
			\left\{
			\inf_{\theta\in\Theta_m}
			\mathrm{KL}\!\big(P^{\star}\,\|\,f_{m,\theta}\big)
			+ \frac{(\log 2)\,L(m)}{n}
			\right\}.
		\end{equation}
	\end{theorem}
	
	\begin{proof}
	Write $\rho(f,g)=\int\sqrt{fg}\,d\mu\in(0,1]$ for the Hellinger affinity,
	so that $h^2(f,g)=1-\rho(f,g)$. Replacing $L$ by $2L$ if necessary (which
	preserves Kraft~\eqref{eq:kraft} and only doubles the penalty constant), we
	may assume the half-exponent condition
	\begin{equation}\label{eq:half-kraft}
		\sum_{m\in\mathcal{I}} 2^{-L(m)/2} \le 1 .
	\end{equation}

	\emph{Step 1 (Kraft/partition-function bound).}
	Since the MCE maximizes $2^{-L(m)}\prod_{i=1}^n f_m(X_i)$ over $m$
	(this is \eqref{eq:def-mce} with code lengths in base $2$), taking square
	roots gives
	\[
	2^{-L(\widehat m)/2}\prod_{i=1}^n\sqrt{\frac{\widehat f(X_i)}{f^{\star}(X_i)}}
	= \max_{m}\, 2^{-L(m)/2}\prod_{i=1}^n\sqrt{\frac{f_m(X_i)}{f^{\star}(X_i)}}
	\le \sum_{m} 2^{-L(m)/2}\prod_{i=1}^n\sqrt{\frac{f_m(X_i)}{f^{\star}(X_i)}}.
	\]
	Taking $\mathbb{E}_{P^{\star}}$ and using independence together with
	$\mathbb{E}_{P^{\star}}\sqrt{f_m/f^{\star}}=\rho(f^{\star},f_m)\le 1$,
	\begin{equation}\label{eq:partition-bound}
		\mathbb{E}_{P^{\star}}\!\left[
		2^{-L(\widehat m)/2}\prod_{i=1}^n\sqrt{\frac{\widehat f(X_i)}{f^{\star}(X_i)}}
		\right]
		\le \sum_{m} 2^{-L(m)/2}\,\rho(f^{\star},f_m)^n
		\le \sum_{m} 2^{-L(m)/2} \le 1 .
	\end{equation}
	This rigorous estimate, a consequence of Kraft's inequality alone, is what
	replaces the (incorrect) claim $\mathbb{E}[L(\widehat m)]\le C_0$; in
	particular it does \emph{not} use any in-sample identity for
	$\mathrm{KL}(P^{\star}\,\|\,\widehat f)$, which would ignore the
	empirical-process fluctuation of Lemma~\ref{lem:KL-decomp}.

	\emph{Step 2 (resolvability bound).}
	Inequality \eqref{eq:partition-bound} is the central estimate of the
	Barron--Cover theory \citep{BarronRissanenYu1998}. Passing from it to the
	index of resolvability is standard: combining \eqref{eq:partition-bound}
	with Jensen's inequality and the elementary comparisons $h^2\le-\log\rho$
	and $-\log\rho(f^{\star},f_m)\le \mathrm{KL}(P^{\star}\,\|\,f_m)$ yields
	\begin{equation}\label{eq:resolvability}
		\mathbb{E}_{P^{\star}}\big[h^2(P^{\star},\widehat f)\big]
		\le \inf_{m\in\mathcal{I}}
		\left\{ \mathrm{KL}(P^{\star}\,\|\,f_m) + \frac{(\log 2)\,L(m)}{n}\right\}.
	\end{equation}
	For a model $\mathcal{M}_m=\{f_{m,\theta}\}$ with continuous $\Theta_m$, the
	same argument applied to a discretization of $\Theta_m$ (whose code length
	is absorbed into $L(m)$) gives \eqref{eq:resolvability} with
	$\inf_{\theta\in\Theta_m}\mathrm{KL}(P^{\star}\,\|\,f_{m,\theta})$ in place
	of $\mathrm{KL}(P^{\star}\,\|\,f_m)$, i.e.\ \eqref{eq:oracle-ineq}.
\end{proof}

\begin{remark}[Why Hellinger and not KL]
	The bound \eqref{eq:oracle-ineq} is stated for the squared Hellinger risk,
	which is bounded by $1$ and is the natural loss in the Barron--Cover theory.
	The analogue with $\mathrm{KL}(P^{\star}\,\|\,\widehat f)$ on the left is
	\emph{false} in general: with positive probability the penalized maximizer
	may select a density $\widehat f$ vanishing faster than $f^{\star}$ in the
	tails, so that $\mathrm{KL}(P^{\star}\,\|\,\widehat f)=+\infty$ while the
	resolvability term stays finite. A KL-risk oracle inequality therefore
	requires a further condition bounding the log-likelihood ratios
	$\log(f^{\star}/f_{m,\theta})$, under which $h^2$ and $\mathrm{KL}$ are
	equivalent up to a constant.
\end{remark}
	
	\subsection{Minimax lower bound}
	
	We now prove a minimax lower bound showing that the oracle rate
	appearing in Theorem~\ref{thm:oracle} is essentially optimal.
	
	\begin{theorem}[Minimax lower bound]
		\label{thm:minimax-lower-proof}
		Let $\mathcal{F}$ be a model class containing a finite subset
		$\{f_0,\ldots,f_M\}$ such that
		\begin{enumerate}
			\item $\mathrm{KL}(f_j\,\|\,f_k) \le \alpha$ for all $j\ne k$;
			\item $\mathrm{KL}(f_j\,\|\,f_0) \ge \beta$ for all $j\ge 1$.
		\end{enumerate}
		Then, for any estimator $\widehat{f}$ based on $n$ observations, we have
		\begin{equation}
			\label{eq:minimax-lb}
			\sup_{f\in\mathcal{F}}
			\mathbb{E}_f\big[\mathrm{KL}(f\,\|\,\widehat{f})\big]
			\ge c\,\beta\,
			\left(1 - \frac{n\alpha + \log 2}{\log M}\right),
		\end{equation}
		for a universal constant $c>0$.
	\end{theorem}
	
	\begin{proof}
		The proof follows the standard Fano--type argument. Consider the
		multi-hypothesis testing problem where the true distribution is
		uniformly drawn from $\{f_0,\ldots,f_M\}$ and one observes
		$X^{(n)}=(X_1,\ldots,X_n)$ i.i.d.\ from the chosen distribution.
		Let $J$ denote the random index of the true distribution and let
		$\widehat{J}$ be any estimator of $J$ based on $X^{(n)}$.
		
		On the one hand, by the data processing inequality and the chain rule
		for Kullback--Leibler divergence,
		\[
		I(J;X^{(n)})
		\le \frac{1}{M+1}\sum_{j=0}^M
		\mathrm{KL}\big(f_j^{\otimes n}\,\|\,\bar{f}^{\otimes n}\big),
		\]
		where $\bar{f} = \frac{1}{M+1}\sum_{j=0}^M f_j$.
		A direct calculation (e.g.\ \citealp{Tsybakov2009}) yields
		\[
		I(J;X^{(n)})
		\le n\alpha,
		\]
		up to negligible terms, under assumption (i).
		
		On the other hand, Fano's inequality states that
		\[
		\mathbb{P}(J\ne\widehat{J})
		\ge 1 - \frac{I(J;X^{(n)})+\log 2}{\log(M+1)}.
		\]
		If $\widehat{f}$ is any estimator, define $\widehat{J}$ by
		\[
		\widehat{J}
		\in \arg\min_{0\le j\le M} \mathrm{KL}(f_j\,\|\,\widehat{f}).
		\]
		Then
		\[
		\sup_{f\in\mathcal{F}}
		\mathbb{E}_f\big[\mathrm{KL}(f\,\|\,\widehat{f})\big]
		\ge \frac{1}{M+1}
		\sum_{j=0}^M
		\mathbb{E}_{f_j}\big[\mathrm{KL}(f_j\,\|\,\widehat{f})\big].
		\]
		Using assumption (ii), whenever $\widehat{J}\ne J$ we have
		$\mathrm{KL}(f_J\,\|\,\widehat{f})\ge \beta/2$ (for a suitable choice
		of the subset). Hence
		\[
		\frac{1}{M+1}
		\sum_{j=0}^M
		\mathbb{E}_{f_j}\big[\mathrm{KL}(f_j\,\|\,\widehat{f})\big]
		\ge \frac{\beta}{2}\,
		\mathbb{P}(J\ne\widehat{J}).
		\]
		Combining the last three displays and simplifying the constants yields
		\eqref{eq:minimax-lb}.
	\end{proof}
	
	\subsection{Consistency of the minimum complexity estimator}
	
	We turn to asymptotic properties.  The following theorem states that,
	under mild identifiability and regularity conditions, the MCE is
	strongly consistent.
	
	\begin{theorem}[Strong consistency]
		\label{thm:consistency-proof}
		Assume that:
		\begin{enumerate}
			\item The true distribution $P^{\star}$ belongs to the union
			$\bigcup_{m\in\mathcal{I}}\mathcal{M}_m$ and is uniquely
			represented as $f_{m^{\star},\theta^{\star}}$;
			\item For each $m$, the parameter space $\Theta_m$ is compact and
			the map $\theta\mapsto f_{m,\theta}$ is continuous in $\mathrm{KL}$;
			\item The penalty satisfies $L(m) = o(n)$ as $n\to\infty$ for each
			fixed $m$ and grows sufficiently fast in $m$ to avoid overfitting.
		\end{enumerate}
		Then
		\begin{equation}
			\widehat{m}\to m^{\star},\qquad
			\widehat{\theta}\to\theta^{\star}
		\end{equation}
		almost surely as $n\to\infty$.
	\end{theorem}
	
	\begin{proof}
		Following the standard M-estimation paradigm, define the normalized
		criterion
		\[
		\mathcal{C}_n(m,\theta)
		:= -\frac{1}{n}\sum_{i=1}^n \ell_{m,\theta}(X_i)
		+ \frac{L(m)}{n}.
		\]
		By the strong law of large numbers and continuity in $\theta$,
		\[
		\mathcal{C}_n(m,\theta)
		\to \mathcal{C}(m,\theta)
		:= - P^{\star}\ell_{m,\theta}
		+ 0
		= \mathrm{KL}(P^{\star}\,\|\,f_{m,\theta}) + \mathrm{const.}
		\]
		almost surely, uniformly on compact subsets of $\Theta_m$.
		By identifiability, $(m^{\star},\theta^{\star})$ is the unique minimizer
		of $\mathcal{C}(m,\theta)$.  The additional growth condition on the
		penalty ensures that for all $m\ne m^{\star}$,
		\[
		\inf_{\theta\in\Theta_m} \mathcal{C}(m,\theta)
		> \mathcal{C}(m^{\star},\theta^{\star}).
		\]
		An application of the argmin continuity theorem (see, e.g.,
		\citealp{vdVart1998}, Theorem 5.7) shows that any sequence of
		approximate minimizers of $\mathcal{C}_n$ converges almost surely
		to $(m^{\star},\theta^{\star})$.  Since $(\widehat{m},\widehat{\theta})$
		exactly minimizes $\mathcal{C}_n$, the result follows.
	\end{proof}
	
	\subsection{Asymptotic distribution}
	
	Under additional regularity (differentiability in quadratic mean,
	non-singular Fisher information, etc.), one can refine the previous
	consistency result into a central limit theorem for $\widehat{\theta}$.
	
	\begin{theorem}[Asymptotic normality]
		\label{thm:asymp-normal-proof}
		Assume, in addition to the hypotheses of Theorem~\ref{thm:consistency-proof},
		that the model $\mathcal{M}_{m^{\star}}$ is regular in the sense of
		classical likelihood theory, with Fisher information matrix
		$I(\theta^{\star})$ positive definite.  Then
		\begin{equation}
			\sqrt{n}\,(\widehat{\theta}-\theta^{\star})
			\xrightarrow{d}
			\mathcal{N}\big(0, I(\theta^{\star})^{-1}\big).
		\end{equation}
	\end{theorem}
	
	\begin{proof}[Sketch of proof]
		We only give the main steps, as they closely follow
		standard arguments in the theory of maximum likelihood estimators.
		Write the score function
		\[
		S_n(\theta)
		:= \frac{1}{\sqrt{n}}
		\sum_{i=1}^n \nabla_{\theta}\ell_{m^{\star},\theta}(X_i)
		\]
		and observe that the contribution of the penalty $L(m)/n$ is negligible
		for fixed $m^{\star}$ (since $L(m^{\star})=o(n)$).
		Expanding the score around $\theta^{\star}$ and using the fact that
		$S_n(\widehat{\theta})\approx 0$ at the optimum, we obtain
		\[
		0
		= S_n(\theta^{\star})
		+ \nabla_{\theta}^2\ell_{m^{\star},\bar{\theta}}(X^{(n)})
		\sqrt{n}\,(\widehat{\theta}-\theta^{\star}),
		\]
		for some intermediate point $\bar{\theta}$, with high probability.
		The central limit theorem applied to $S_n(\theta^{\star})$, together
		with a law of large numbers for the observed information matrix,
		implies that
		\[
		S_n(\theta^{\star})
		\xrightarrow{d} \mathcal{N}(0,I(\theta^{\star})),\qquad
		-\frac{1}{n}
		\nabla_{\theta}^2\ell_{m^{\star},\bar{\theta}}(X^{(n)})
		\xrightarrow{\mathbb{P}} I(\theta^{\star}),
		\]
		and the result follows by Slutsky's lemma.
	\end{proof}
	
	\subsection{Further auxiliary lemmas}
	
	For completeness, we collect here two additional lemmas on covering
	numbers and metric entropy that are used in the detailed nonasymptotic
	bounds of Section~\ref{sec:nonasymptotic}.
	
	\begin{lemma}[Covering numbers]
		\label{lem:covering}
		Let $\mathcal{F}$ be a class of functions with envelope $F$ and let
		$N(\varepsilon,\mathcal{F},L_2(P^{\star}))$ denote its covering number
		with respect to the $L_2(P^{\star})$--metric.  If there exist constants
		$A,v>0$ such that
		\[
		\log N(\varepsilon,\mathcal{F},L_2(P^{\star}))
		\le v \log\left(\frac{A}{\varepsilon}\right)
		\]
		for all $0<\varepsilon\le A$, then the empirical process indexed by
		$\mathcal{F}$ satisfies
		\[
		\mathbb{E}
		\sup_{f\in\mathcal{F}}
		\big|(P_n-P^{\star})f\big|
		\le C \sqrt{\frac{v}{n}}\,
		\|F\|_{L_2(P^{\star})},
		\]
		for a universal constant $C>0$.
	\end{lemma}
	
	\begin{proof}
		This is a standard consequence of Dudley's entropy integral and can
		be found in many references on empirical process theory; see, e.g.,
		\cite{vdVW1996}, Theorem 2.14.2.
	\end{proof}
	
	\begin{lemma}[Local entropy bound]
		\label{lem:local-entropy}
		Under the same hypotheses as Lemma~\ref{lem:covering}, the local
		entropy integral on a ball of radius $r$ in $L_2(P^{\star})$ is bounded
		by
		\[
		\int_0^r
		\sqrt{\log N(\varepsilon,\mathcal{F},L_2(P^{\star}))}\,d\varepsilon
		\le C' r \sqrt{v \log\left(\frac{A}{r}\right)},
		\]
		for a universal constant $C'>0$.
	\end{lemma}
	
	\begin{proof}
		Integrate the assumed entropy bound and use Jensen's inequality.
	\end{proof}
	
	These lemmas allow one to replace the crude union bound in
	Lemma~\ref{lem:exp-ineq} by sharper chaining arguments when needed.
	We have opted to present the main ideas in a form that keeps the
	structure of the proofs transparent while remaining compatible with
	the level of rigor expected in leading probability and statistics
	journals such as \emph{The Annals of Statistics}.
	
	\bigskip
	
	This completes the proofs of the main results.
	
	\begin{thebibliography}{99}
		
		\bibitem[ Barron, Rissanen and Yu(1998)]{BarronRissanenYu1998}
		Barron, A., Rissanen, J. and Yu, B. (1998).
		The minimum description length principle in coding and modeling.
		\emph{IEEE Trans. Inform. Theory}, \textbf{44}(6), 2743--2760.

		\bibitem[Birg\'e and Massart(1998)]{BirgeMassart1998}
		Birg\'e, L. and Massart, P. (1998).
		Minimum contrast estimators on sieves: exponential bounds and rates of
		convergence. \emph{Bernoulli}, \textbf{4}(3), 329--375.

		\bibitem[Massart(2007)]{Massart2007}
		Massart, P. (2007).
		\emph{Concentration Inequalities and Model Selection}.
		Lecture Notes in Mathematics 1896, Springer, Berlin.
		
		\bibitem[Tsybakov(2009)]{Tsybakov2009}
		Tsybakov, A. (2009).
		\emph{Introduction to Nonparametric Estimation}.
		Springer, New York.
		
		\bibitem[van der Vaart and Wellner(1996)]{vdVW1996}
		van der Vaart, A. W. and Wellner, J. A. (1996).
		\emph{Weak Convergence and Empirical Processes}.
		Springer, New York.
		
		\bibitem[van der Vaart(1998)]{vdVart1998}
		van der Vaart, A. W. (1998).
		\emph{Asymptotic Statistics}.
		Cambridge University Press, Cambridge.
		
	\end{thebibliography}
	
\end{document}
